\section{Recovery that preserves a causal request}
\label{exec14:failure-trials}

A recoverable service must preserve the scientific request as well as the
computation. Suppose a servicing analyst requests an effect in dollars over
twelve months and the calculation is interrupted. A later response for euros,
another horizon or a revised graph is not a continuation of that request.
Conversely, the absence of a response says nothing about whether the effect is
zero, unidentifiable or economically unattractive. Our proposed integration
therefore maintains separate identities for the scientific context, execution
environment, input artifacts, completed work and published result.
\citep{exec14_failure_design}

The local reference runtime already binds requests to context and artifact
identities, and binds replay to selected code, schemas, Python and dependency
versions. Its report jobs preserve a prefix of completed artifact identifiers
and corresponding summaries. Reports describe supplied results as declarations
checked against a schema; they do not confer independent causal review. These
are useful mechanics for an early prototype, with a deliberately narrow
single-operator scope. They neither authenticate the supplied business premises
nor qualify the system for bank operation. \citep{exec14_failure_implementation}

\subsection{Semantic failure and operational failure require different responses}

Generated local requests illustrate the distinctions. A valid quadratic
sensitivity request returned a computed empty set. Changing its declared units
from USD to EUR while retaining the original input artifact produced a failed,
unavailable response with no mathematical result. An unsupported discovery
method returned a completed invocation with the scientific disposition
unsupported and no result. A one-byte output limit produced an output-budget
failure. These observations demonstrate distinct response meanings for selected
requests; an empty mathematical set must not be used as a generic representation
of missing computation. \citep{exec14_failure_observations}

The review also exposed a boundary that the proposed interface must repair.
The library raises a validation exception for an unsupported request-schema
version. It also raises a JSON decoding exception when a previously saved replay
file is manually truncated to an incomplete JSON object. The latter case used
an unchanged, valid request; the fault was in local saved state. A follow-up
through the actual CLI found that both conditions return exit code two and
the error code invalid\_request. The CLI catches the exception, but assigns
the storage fault to the wrong operational category. A caller could repeatedly
rewrite a scientifically valid request while the corrupted replay state remains
unchanged. \citep{exec14_failure_observations}

This finding does not establish that the runtime's atomic writer naturally emits
truncated files. The truncation was a deliberate local fault injection. It does
establish the observed reader behavior for those bytes, and a concrete recovery
requirement. The proposed response should classify invalid caller syntax
separately from corrupt or unavailable saved state, preserve the offending bytes
and their identity for diagnosis, and expose whether retry can succeed. It should
not silently delete the cache and rerun a scientific operation whose original
inputs, environment, randomization or authorization may no longer be available.
A rebuild must first revalidate those identities and declare whether it resumes
the old request or creates a new attempt. \citep{exec14_failure_design}

\subsection{Three interruption windows}

The content-addressed artifact store publishes bytes and metadata through
separate atomic writes. Individual write atomicity does not make the pair one
transaction. In a generated case we paused immediately after publishing the
content file, before publishing its descriptor, and killed that child process.
The resulting content-only artifact was absent from the reference listing and
could not be retrieved by descriptor lookup. Retrying the same value, schema
and role completed the missing metadata under the same handle and returned the
original value. This supports one idempotent repair route. A different value
under that handle would require a different decision; this observation does not
authorize silently reusing a stale descriptor. \citep{exec14_failure_observations}

The second window concerns work already completed but not yet delivered.
A report over three generated canonical responses was interrupted after its
first saved summary. An unmodified runtime then resumed the job. The final
ordered identifiers exactly matched the original three requested artifacts,
and the report contained one summary per identifier. Every summary retained
the qualification that checking was schema-only, without causal or external
review. Preserving that qualification across recovery is as necessary as
avoiding duplicated work: a restart must not upgrade a conditional result into
an accepted finding. \citep{exec14_failure_observations}

The third window occurs after the report artifact has been stored but before the
job's terminal state points to it. We paused at that exact boundary and killed
the child. Saved job state still reported running with no result reference,
although the report artifact already existed. Recovery completed the job and
returned the same report handle and bytes. Thus a durable content-addressed
artifact can support idempotent publication even when acknowledgement is lost.
The demonstrated behavior concerns local deterministic summaries. A future
external side effect, such as submitting a bank instruction, would need its own
transaction and reconciliation protocol; repeating a local report write supplies
no evidence about that different operation. \citep{exec14_failure_observations,exec14_failure_design}

These three trials used explicit test-process hooks at named publication
boundaries. The library source files were unchanged. The test child executed
the actual store or report-worker code up to the boundary, and the parent killed
only its own child. Subsequent recovery used the ordinary runtime. Initial
submission scheduling was deliberately deferred so the selected boundary could
be reached deterministically; therefore the production spawn counter excludes
that injected first worker. We do not treat the displayed counter as a count
of every process launched by the harness. These cases give reproducible
mechanism evidence, not outage frequencies or coverage of every operating-system
failure ordering. \citep{exec14_failure_observations}

\subsection{A checkpoint is evidence to validate, not an instruction to trust}

In another generated case, the first saved summary was changed while retaining
its old checkpoint digest. Resume failed with a checkpoint-content mismatch,
kept the result reference absent and preserved the failed state. This is the
appropriate local refusal for that mutation. The implementation also checks
that completed identifiers form a prefix of the frozen request, that summary
and identifier counts agree, and that execution identity remains compatible.
These checks bind a particular continuation. They do not authenticate a writer
who can coherently replace every relevant descriptor and trust anchor; the
separate custody discussion retains that limitation. \citep{exec14_failure_implementation,exec14_failure_observations}

The proposed state transition can be expressed without assuming success:
\[
 \text{queued}\longrightarrow\text{running}
 \longrightarrow\{\text{completed},\text{failed},\text{cancelled}\}.
\]
Resumption of unfinished work is allowed only after the input, execution and
checkpoint conditions hold. Completion additionally requires a readable result
whose content and meaning match the declared request. Cancellation records a
terminal decision about future work; it is not a claim that previously published
artifacts or external effects have been rolled back. A nonterminal job may
already have a durable artifact, as the third interruption case demonstrates.
The authoritative job state and the artifact's independent identity must remain
distinguishable. \citep{exec14_failure_design}

For the next prototype, the known repair is to centralize error normalization
for library, CLI and MCP callers, while preserving the difference between
request defects, state corruption, resource exhaustion and scientific refusal.
Acceptance cases should include malformed and missing replay records, wrong
context or schema, unreadable referenced artifacts, changed code identity,
partial metadata, corrupt checkpoints, cancellation at each publication boundary,
and recovery near the original deadline. Negative cases must remain in the
request ledger with their attempted inputs and observed disposition. A valid
retry must not gain a fresh unlimited budget merely by restarting.
\citep{exec14_failure_design}

The current observations establish selected local behavior and expose a concrete
classification defect. They do not repair the implementation, rerun a live MCP
consumer, measure bank-scale recovery, or settle the previously observed
descriptor-replacement and FIFO limitations. Those claims require their own
tests and, for operating acceptance, a named consumer, target filesystem and
platform, identity authority, retention rules and declared service limits.
Public or generated fixtures can justify the prototype and demonstrate the
merit of later experiments before expensive bank trials are commissioned.
\citep{exec14_failure_design}
\label{exec14:failure-trials-end}
